\section{Conclusion}
\label{sec:conclusion}

TRACE-BiOpt reframes sparse traffic sensor deployment as a recoverability-driven bilevel network design problem. Instead of ranking locations by coverage, variance, or topology alone, it selects a fixed sensor set by directly optimizing hidden-network reconstruction risk under a transparent GLS/MAP estimator.

That is also the comparison class the paper asks a reviewer to use: fixed-infrastructure transportation network design under downstream recoverability, not estimator tuning, benchmark-style model ranking, or post-hoc layout selection.

The method combines hidden Huber reconstruction loss, posterior uncertainty, scenario tail risk, and spatial regularization in one objective, then applies deterministic initialization and exchange refinement under that objective. Across PeMS7\_228, PeMS7\_1026, and Seattle current-best evidence, TRACE-BiOpt beats the best pre-registered non-BiOpt baseline in mean held-out GLS/MAP MAE in every tested dataset-budget regime, the audited comparison class spans 21 baselines across 11 method families, the current nine-row chain is submission-ready under paired audits, Holm-corrected one-sided paired tests across all 189 current-best row-baseline comparisons leave no statistically tied or significantly better pre-registered challenger, and the bounded exact benchmark exact-hits 27/27 deterministic 16-node subnetworks with zero objective gap. The final contribution is an auditable bridge between traffic infrastructure design, transparent inverse problems, robust bilevel stochastic optimization, and empirical sensor deployment evidence.

Stated in the most compressed TR-B form, the paper claims one thing: one
fixed-infrastructure transportation network-design method built around a
transparent inverse problem can be made formulation-consistent, theory-backed,
compressible to row-wise strongest-baseline dominance, multiplicity-robust
against the full audited comparison class, and interpretable as a principled
bilevel stochastic design method on the tested traffic networks without
turning into a post-hoc selector story.

Table~\ref{tab:trace-biopt-planning-takeaways} turns that bridge into a
network-level planning readout. The point is not to add another dominance
table. The point is to close with what a transport-planning reader can
actually do with the evidence: treat PeMS7\_1026 as a search-budget-sensitive
large network, PeMS7\_228 as an expansion-friendly calibrated rollout, and
Seattle as a conservative staged-rollout case whose 20/30\% promoted routes
are deployable while the 10\% row remains intentionally fail-closed on
promotion evidence.

\begin{table*}[t]
\centering
\caption{Network-level planning takeaways from the TRACE-BiOpt current-best evidence chain. Rather than restating dominance, the table translates each tested network into a transport-planning posture for staged deployment decisions.}
\label{tab:trace-biopt-planning-takeaways}
\scriptsize
\begin{tabularx}{\textwidth}{>{\raggedright\arraybackslash}p{0.13\textwidth}>{\raggedright\arraybackslash}p{0.20\textwidth}>{\raggedright\arraybackslash}p{0.31\textwidth}>{\raggedright\arraybackslash}X}
\toprule
Dataset & Planning archetype & Current-best readout & Transport-planning takeaway \\
\midrule
PeMS7\_1026 & Large search-sensitive network & All three budgets are Stage16-promoted; strongest-baseline margins contract 0.130 -> 0.124 -> 0.038 even as TRACE-BiOpt MAE keeps falling, the Holm-corrected within-dataset all-baseline screen stays clean at 63/63, and the bounded exact benchmark still hits 9/9 audited subnetworks. & Treat added sensors and added search budget as complementary capital: larger networks need calibrated risk estimation and deeper deterministic exchange search before a weak margin should be interpreted as method failure. \\
PeMS7\_228 & Expansion-friendly medium network & All three budgets are Stage16-promoted; strongest-baseline margins widen 0.142 -> 0.276 -> 0.336, the Holm-corrected within-dataset all-baseline screen stays clean at 63/63, and the bounded exact benchmark stays at 9/9 audited hits. & Once low-budget calibration is repaired, staged expansion is justified: extra hardware keeps paying off in both absolute reconstruction error and relative margin over the strongest audited comparator. \\
Seattle & Low-variance staged-rollout network & The 20\% and 30\% rows are Stage16-promoted while Stage15 main evidence remains for Seattle 10\%; the strongest-baseline margins still widen 0.056 -> 0.083 -> 0.095, the Holm-corrected within-dataset all-baseline screen stays clean at 63/63, and bounded exact subnetworks stay at 9/9 hits without upgrading the fail-closed 10\% promotion lane. & Seattle is the cleanest staged-rollout case, but the conservative reading should remain visible: keep the 10\% row on its retained evidence lane and use the promoted 20/30\% routes as the operationally preferred default. \\
\bottomrule
\end{tabularx}
\end{table*}


In that sense, the paper's claim is not just that one siting heuristic wins a
benchmark. It is that a recoverability-driven network-design formulation, a
single-objective deterministic solver, a scoped but explicit theory package,
and an external audited comparison-class current-best evidence chain can be
made to line up in one transport-network design argument without hiding behind
a candidate pool.
The intended reading is therefore that TRACE-BiOpt should be evaluated as one
fixed-infrastructure transportation network-design method built around a
transparent inverse problem and judged under an external audited
comparison-class contract that compresses to row-wise strongest-baseline
dominance on held-out reconstruction error and then survives all-baseline
corrected paired tests, not as a sparse-sensor benchmark paper that happens to
borrow traffic data.
Its shortest reviewer-facing summary is the same one used at the front of the
paper: one long-lived infrastructure decision, one transparent inverse
problem, one deterministic solver, one scoped theory package, and one
external audited comparison-class contract over 21 pre-registered baselines
spanning 11 method families with no surviving challenger after 189
Holm-corrected paired tests.
For a transportation agency, the practical message is equally narrow and
concrete: freeze the sparse sensor siting plan as an infrastructure design decision
before deployment, evaluate it by the recoverability it creates for the hidden
network state under a transparent estimator, and compare it to the strongest
audited alternative rather than to a weak fixed baseline or a post-hoc
estimator gain. That is the decision contract this paper argues for.
